\section{Definitions and tools}\label{sec:setup}

\subsection{Blackouts}

An $h\times m$ grid has $h$ rows, $m$ columns, and point set $V$. The point
$(i,j)$ lies in column $i$ and row $j$, with $(1,1)$ at the top left
(Figure~\ref{fig:game}(a)). A \emph{rectangle} has four grid points as corners
and positive area. We write $C(R)$ for the corner set of a rectangle $R$. It is
\emph{axis-aligned} if its sides follow the grid lines, and \emph{tilted}
otherwise.

A \emph{blackout} is a set $S\subseteq V$ of black points. The other points are
white. The \emph{presentation} of a rectangle $R$ is $\pres(R)=C(R)\setminus S$,
the set of its white corners. We say that $R$ \emph{shows} these corners, and
that $R$ is \emph{hidden} if it shows none.

\begin{definition}[valid blackout]\label{def:valid}
A blackout $S$ is \emph{valid} if distinct rectangles have distinct
presentations. We write $f(h,m)$ for the largest size of a valid blackout
of the $h\times m$ grid, and $a_h(m)$ for the number of valid blackouts of that
size, which we call \emph{optima}. We count them as sets of points, so
reflections and rotations are not identified.
\end{definition}

Transposing the grid gives $f(h,m)=f(m,h)$ and $a_h(m)=a_m(h)$. A grid with
$1$ row or $1$ column has no rectangles, so $f(1,m)=m$ and $a_1(m)=1$.

Column $i$ of a blackout has \emph{black rows} $B_i=\{j:(i,j)\in S\}$. It is
\emph{empty}, a \emph{singleton}, a \emph{double}, or a \emph{triple} when
$|B_i|$ is $0$, $1$, $2$, or $3$, and \emph{full} when $|B_i|=h$. The \emph{column
word} of $S$ is $B_1B_2\cdots B_m$. The \emph{pair set} $F$ of $S$ is the set of
pairs of rows that are both black in some column. Throughout,
$q_h=\lfloor h^2/4\rfloor$, which is \seqnum{A002620}$(h)$.

\begin{lemma}[difference sets]\label{lem:diff}
Two rectangles $R_1$ and $R_2$ have the same presentation if and only if
$C(R_1)\triangle C(R_2)\subseteq S$.
\end{lemma}

\begin{proof}
The sets $C(R_1)\setminus S$ and $C(R_2)\setminus S$ agree exactly when every
point that lies in just one of $C(R_1)$ and $C(R_2)$ belongs to $S$.
\end{proof}

\begin{lemma}[three corners suffice]\label{lem:threecorners}
Three corners of a rectangle determine it. Hence two distinct rectangles with the
same presentation show at most $2$ corners, and a valid blackout hides at most
$1$ rectangle.
\end{lemma}

\begin{proof}
Three corners form a right triangle. Its right angle is at the corner $q$
adjacent to the other two, $p$ and $r$, so the fourth corner is $p+r-q$. Two
hidden rectangles would both have the empty presentation.
\end{proof}

\subsection{Tilted rectangles}

Put $W_h=1+\lfloor(h-1)^2/4\rfloor$, so $W_3=2$, $W_4=3$, and $W_5=5$.

\begin{lemma}[leftmost corner]\label{lem:leftmost}
A tilted rectangle $R$ has a unique corner $\ell=(x,r)$ in its smallest column.
The sides from $\ell$ are $(a,-b)$ and $(c,d)$ with positive integers $a,b,c,d$
and $ac=bd$, and the corners are
\begin{equation}\label{eq:tilt-param}
(x,r),\qquad(x+a,\,r-b),\qquad(x+c,\,r+d),\qquad(x+a+c,\,r-b+d).
\end{equation}
The rectangle meets the rows $r-b$ to $r+d$ and the columns $x$ to $x+a+c$. We
call $a+c$ its \emph{span}.
\end{lemma}

\begin{proof}
No side of $R$ is vertical. Hence two corners in the smallest column $x$ would
be opposite. The other two corners have the same midpoint, so they would also
lie in column $x$, and $R$ would have no area. The two sides from
$\ell$ therefore have positive first coordinates $a$ and $c$. They are
perpendicular, so their second coordinates have product $-ac$. Thus one
side points up and the other down, which gives the form above with $ac=bd$.
\end{proof}

\begin{lemma}[span bound]\label{lem:span}
Every tilted rectangle of the $h\times m$ grid has span $a+c\le W_h$.
\end{lemma}

\begin{proof}
The rectangle fits in $h$ rows, so $b+d\le h-1$. Since $(a-1)(c-1)\ge0$, we get
$a+c\le ac+1=bd+1$. Finally $bd\le\lfloor(b+d)^2/4\rfloor\le W_h-1$.
\end{proof}

\begin{remark}[a conjecture in \seqnum{A338885}]\label{cor:oeisconj}
The same computation settles a conjecture of Kagey. Row $n$ of his triangle
\seqnum{A338885} lists all $k$ such that some tilted lattice rectangle touches
all $4$ sides of an $n\times k$ rectangle, and he conjectures there that the
largest such $k$ is \seqnum{A033638}$(n)=\lfloor n^2/4\rfloor+1$. This is true
for every $n\ge2$. In~\eqref{eq:tilt-param} the bounding box has height $b+d=n$
and width $k=a+c$. As in Lemma~\ref{lem:span}, $k\le bd+1\le\lfloor n^2/4\rfloor+1$.
For equality take $b=\lfloor n/2\rfloor$, $d=\lceil n/2\rceil$, and the
perpendicular sides $(1,-b)$ and $(bd,d)$.
\end{remark}

The \emph{diamond} $D(i,j)$ is the square with corners $(i,j+1)$, $(i+1,j)$,
$(i+2,j+1)$, and $(i+1,j+2)$. The rectangles $\kA_i$, $\kB_i$, $\kC_i$, and
$\kE_i$ have one corner in each of the columns $i$ to $i+3$, in the rows
$3412$, $3142$, $2413$, and $2143$ from left to right (Figure~\ref{fig:tilted}).

\begin{figure}[t]
\centering
\begin{tikzpicture}[x=0.5cm,y=0.5cm,baseline=(current bounding box.north)]
  \begin{scope}[xshift=0.000cm]
    \draw[black,line width=0.20pt,dash pattern=on 0.75pt off 0.75pt] (1,-1) -- (4,-1) (1,-2) -- (4,-2) (1,-3) -- (4,-3) (1,-4) -- (4,-4) (1,-1) -- (1,-4) (2,-1) -- (2,-4) (3,-1) -- (3,-4) (4,-1) -- (4,-4);
    \draw[kred,line width=1.42pt,line join=round] (1,-2) -- (2,-1) -- (3,-2) -- (2,-3) -- cycle;
    \filldraw[fill=white,draw=black,line width=0.30pt] (1,-1) circle (0.3);
    \filldraw[fill=white,draw=black,line width=0.30pt] (2,-1) circle (0.3);
    \filldraw[fill=white,draw=black,line width=0.30pt] (3,-1) circle (0.3);
    \filldraw[fill=white,draw=black,line width=0.30pt] (4,-1) circle (0.3);
    \filldraw[fill=kyellow,draw=black,line width=0.30pt] (1,-2) circle (0.3);
    \filldraw[fill=white,draw=black,line width=0.30pt] (2,-2) circle (0.3);
    \filldraw[fill=white,draw=black,line width=0.30pt] (3,-2) circle (0.3);
    \filldraw[fill=white,draw=black,line width=0.30pt] (4,-2) circle (0.3);
    \filldraw[fill=white,draw=black,line width=0.30pt] (1,-3) circle (0.3);
    \filldraw[fill=white,draw=black,line width=0.30pt] (2,-3) circle (0.3);
    \filldraw[fill=white,draw=black,line width=0.30pt] (3,-3) circle (0.3);
    \filldraw[fill=white,draw=black,line width=0.30pt] (4,-3) circle (0.3);
    \filldraw[fill=white,draw=black,line width=0.30pt] (1,-4) circle (0.3);
    \filldraw[fill=white,draw=black,line width=0.30pt] (2,-4) circle (0.3);
    \filldraw[fill=white,draw=black,line width=0.30pt] (3,-4) circle (0.3);
    \filldraw[fill=white,draw=black,line width=0.30pt] (4,-4) circle (0.3);
    \node[lab] at (1,-0.45) {1};
    \node[lab] at (2,-0.45) {2};
    \node[lab] at (3,-0.45) {3};
    \node[lab] at (4,-0.45) {4};
    \node[lab] at (0.45,-1) {1};
    \node[lab] at (0.45,-2) {2};
    \node[lab] at (0.45,-3) {3};
    \node[lab] at (0.45,-4) {4};
    \node[capt,anchor=north] at (2.50,-4.45) {$D(1,1)$};
  \end{scope}
  \begin{scope}[xshift=2.700cm]
    \draw[black,line width=0.20pt,dash pattern=on 0.75pt off 0.75pt] (1,-1) -- (4,-1) (1,-2) -- (4,-2) (1,-3) -- (4,-3) (1,-4) -- (4,-4) (1,-1) -- (1,-4) (2,-1) -- (2,-4) (3,-1) -- (3,-4) (4,-1) -- (4,-4);
    \draw[kred,line width=1.42pt,line join=round] (1,-3) -- (2,-4) -- (4,-2) -- (3,-1) -- cycle;
    \filldraw[fill=white,draw=black,line width=0.30pt] (1,-1) circle (0.3);
    \filldraw[fill=white,draw=black,line width=0.30pt] (2,-1) circle (0.3);
    \filldraw[fill=white,draw=black,line width=0.30pt] (3,-1) circle (0.3);
    \filldraw[fill=white,draw=black,line width=0.30pt] (4,-1) circle (0.3);
    \filldraw[fill=white,draw=black,line width=0.30pt] (1,-2) circle (0.3);
    \filldraw[fill=white,draw=black,line width=0.30pt] (2,-2) circle (0.3);
    \filldraw[fill=white,draw=black,line width=0.30pt] (3,-2) circle (0.3);
    \filldraw[fill=white,draw=black,line width=0.30pt] (4,-2) circle (0.3);
    \filldraw[fill=kyellow,draw=black,line width=0.30pt] (1,-3) circle (0.3);
    \filldraw[fill=white,draw=black,line width=0.30pt] (2,-3) circle (0.3);
    \filldraw[fill=white,draw=black,line width=0.30pt] (3,-3) circle (0.3);
    \filldraw[fill=white,draw=black,line width=0.30pt] (4,-3) circle (0.3);
    \filldraw[fill=white,draw=black,line width=0.30pt] (1,-4) circle (0.3);
    \filldraw[fill=white,draw=black,line width=0.30pt] (2,-4) circle (0.3);
    \filldraw[fill=white,draw=black,line width=0.30pt] (3,-4) circle (0.3);
    \filldraw[fill=white,draw=black,line width=0.30pt] (4,-4) circle (0.3);
    \node[capt,anchor=north] at (2.50,-4.45) {$\kA_1$};
  \end{scope}
  \begin{scope}[xshift=5.400cm]
    \draw[black,line width=0.20pt,dash pattern=on 0.75pt off 0.75pt] (1,-1) -- (4,-1) (1,-2) -- (4,-2) (1,-3) -- (4,-3) (1,-4) -- (4,-4) (1,-1) -- (1,-4) (2,-1) -- (2,-4) (3,-1) -- (3,-4) (4,-1) -- (4,-4);
    \draw[kred,line width=1.42pt,line join=round] (1,-3) -- (3,-4) -- (4,-2) -- (2,-1) -- cycle;
    \filldraw[fill=white,draw=black,line width=0.30pt] (1,-1) circle (0.3);
    \filldraw[fill=white,draw=black,line width=0.30pt] (2,-1) circle (0.3);
    \filldraw[fill=white,draw=black,line width=0.30pt] (3,-1) circle (0.3);
    \filldraw[fill=white,draw=black,line width=0.30pt] (4,-1) circle (0.3);
    \filldraw[fill=white,draw=black,line width=0.30pt] (1,-2) circle (0.3);
    \filldraw[fill=white,draw=black,line width=0.30pt] (2,-2) circle (0.3);
    \filldraw[fill=white,draw=black,line width=0.30pt] (3,-2) circle (0.3);
    \filldraw[fill=white,draw=black,line width=0.30pt] (4,-2) circle (0.3);
    \filldraw[fill=kyellow,draw=black,line width=0.30pt] (1,-3) circle (0.3);
    \filldraw[fill=white,draw=black,line width=0.30pt] (2,-3) circle (0.3);
    \filldraw[fill=white,draw=black,line width=0.30pt] (3,-3) circle (0.3);
    \filldraw[fill=white,draw=black,line width=0.30pt] (4,-3) circle (0.3);
    \filldraw[fill=white,draw=black,line width=0.30pt] (1,-4) circle (0.3);
    \filldraw[fill=white,draw=black,line width=0.30pt] (2,-4) circle (0.3);
    \filldraw[fill=white,draw=black,line width=0.30pt] (3,-4) circle (0.3);
    \filldraw[fill=white,draw=black,line width=0.30pt] (4,-4) circle (0.3);
    \node[capt,anchor=north] at (2.50,-4.45) {$\kB_1$};
  \end{scope}
  \begin{scope}[xshift=8.100cm]
    \draw[black,line width=0.20pt,dash pattern=on 0.75pt off 0.75pt] (1,-1) -- (4,-1) (1,-2) -- (4,-2) (1,-3) -- (4,-3) (1,-4) -- (4,-4) (1,-1) -- (1,-4) (2,-1) -- (2,-4) (3,-1) -- (3,-4) (4,-1) -- (4,-4);
    \draw[kred,line width=1.42pt,line join=round] (1,-2) -- (2,-4) -- (4,-3) -- (3,-1) -- cycle;
    \filldraw[fill=white,draw=black,line width=0.30pt] (1,-1) circle (0.3);
    \filldraw[fill=white,draw=black,line width=0.30pt] (2,-1) circle (0.3);
    \filldraw[fill=white,draw=black,line width=0.30pt] (3,-1) circle (0.3);
    \filldraw[fill=white,draw=black,line width=0.30pt] (4,-1) circle (0.3);
    \filldraw[fill=kyellow,draw=black,line width=0.30pt] (1,-2) circle (0.3);
    \filldraw[fill=white,draw=black,line width=0.30pt] (2,-2) circle (0.3);
    \filldraw[fill=white,draw=black,line width=0.30pt] (3,-2) circle (0.3);
    \filldraw[fill=white,draw=black,line width=0.30pt] (4,-2) circle (0.3);
    \filldraw[fill=white,draw=black,line width=0.30pt] (1,-3) circle (0.3);
    \filldraw[fill=white,draw=black,line width=0.30pt] (2,-3) circle (0.3);
    \filldraw[fill=white,draw=black,line width=0.30pt] (3,-3) circle (0.3);
    \filldraw[fill=white,draw=black,line width=0.30pt] (4,-3) circle (0.3);
    \filldraw[fill=white,draw=black,line width=0.30pt] (1,-4) circle (0.3);
    \filldraw[fill=white,draw=black,line width=0.30pt] (2,-4) circle (0.3);
    \filldraw[fill=white,draw=black,line width=0.30pt] (3,-4) circle (0.3);
    \filldraw[fill=white,draw=black,line width=0.30pt] (4,-4) circle (0.3);
    \node[capt,anchor=north] at (2.50,-4.45) {$\kC_1$};
  \end{scope}
  \begin{scope}[xshift=10.800cm]
    \draw[black,line width=0.20pt,dash pattern=on 0.75pt off 0.75pt] (1,-1) -- (4,-1) (1,-2) -- (4,-2) (1,-3) -- (4,-3) (1,-4) -- (4,-4) (1,-1) -- (1,-4) (2,-1) -- (2,-4) (3,-1) -- (3,-4) (4,-1) -- (4,-4);
    \draw[kred,line width=1.42pt,line join=round] (1,-2) -- (3,-4) -- (4,-3) -- (2,-1) -- cycle;
    \filldraw[fill=white,draw=black,line width=0.30pt] (1,-1) circle (0.3);
    \filldraw[fill=white,draw=black,line width=0.30pt] (2,-1) circle (0.3);
    \filldraw[fill=white,draw=black,line width=0.30pt] (3,-1) circle (0.3);
    \filldraw[fill=white,draw=black,line width=0.30pt] (4,-1) circle (0.3);
    \filldraw[fill=kyellow,draw=black,line width=0.30pt] (1,-2) circle (0.3);
    \filldraw[fill=white,draw=black,line width=0.30pt] (2,-2) circle (0.3);
    \filldraw[fill=white,draw=black,line width=0.30pt] (3,-2) circle (0.3);
    \filldraw[fill=white,draw=black,line width=0.30pt] (4,-2) circle (0.3);
    \filldraw[fill=white,draw=black,line width=0.30pt] (1,-3) circle (0.3);
    \filldraw[fill=white,draw=black,line width=0.30pt] (2,-3) circle (0.3);
    \filldraw[fill=white,draw=black,line width=0.30pt] (3,-3) circle (0.3);
    \filldraw[fill=white,draw=black,line width=0.30pt] (4,-3) circle (0.3);
    \filldraw[fill=white,draw=black,line width=0.30pt] (1,-4) circle (0.3);
    \filldraw[fill=white,draw=black,line width=0.30pt] (2,-4) circle (0.3);
    \filldraw[fill=white,draw=black,line width=0.30pt] (3,-4) circle (0.3);
    \filldraw[fill=white,draw=black,line width=0.30pt] (4,-4) circle (0.3);
    \node[capt,anchor=north] at (2.50,-4.45) {$\kE_1$};
  \end{scope}
\end{tikzpicture}

\caption{The tilted rectangles of the $4\times m$ grid, up to translation. Each is
drawn in a $4\times4$ window, and its leftmost corner $\ell$ is yellow.}
\label{fig:tilted}
\end{figure}

\begin{lemma}[tilted rectangles on $3$ and $4$ rows]\label{lem:tilted-narrow}
The tilted rectangles of the $3\times m$ grid are the diamonds $D(i,1)$. Those of
the $4\times m$ grid are $D(i,1)$, $D(i,2)$, $\kA_i$, $\kB_i$, $\kC_i$, and
$\kE_i$.
\end{lemma}

\begin{proof}
We use Lemma~\ref{lem:leftmost}. On $3$ rows, $b+d\le2$ forces $b=d=1$. Then
$ac=1$, so $a=c=1$ and $\ell$ is in row $2$. On $4$ rows, $b=d=1$ again gives
diamonds, with $\ell$ in row $2$ or $3$. Otherwise $\{b,d\}=\{1,2\}$. Then
$ac=2$, so $\{a,c\}=\{1,2\}$, and $\ell$ is in row $1+b$. The $4$ choices of
sides, $(2,-2),(1,1)$, $(1,-2),(2,1)$, $(2,-1),(1,2)$, and $(1,-1),(2,2)$,
give $\kA_i$, $\kB_i$, $\kC_i$, and $\kE_i$.
\end{proof}

\subsection{Axis-aligned rectangles}

A blackout $S$ is \emph{axis-valid} if distinct axis-aligned rectangles have
distinct presentations. In the \emph{row-column graph} $G_S$, the vertices are
the rows $r_1,\dots,r_h$ and the columns $c_1,\dots,c_m$, and $r_j$ is joined to
$c_i$ for each black point $(i,j)$.

\begin{lemma}[short cycles]\label{lem:girth}
If $G_S$ has no $4$-cycle and no $6$-cycle, then $S$ is axis-valid. The converse
holds for $m\ge3$.
\end{lemma}

\begin{proof}
Let $m\ge3$. A $4$-cycle $r_j\,c_i\,r_{j'}\,c_{i'}$ gives the rectangles
$\{i,w\}\times\{j,j'\}$ and $\{i',w\}\times\{j,j'\}$ for a third column $w$.
Their corner sets differ in $4$ black points. A $6$-cycle
$r_a\,c_x\,r_b\,c_y\,r_c\,c_z$ gives $\{x,y\}\times\{a,b\}$ and
$\{y,z\}\times\{a,c\}$, whose corner sets differ in the $6$ black points of the
cycle. In both cases Lemma~\ref{lem:diff} shows that $S$ is not axis-valid.

Conversely, let axis-aligned rectangles $R_1\ne R_2$ have the same presentation,
so that $C(R_1)\triangle C(R_2)\subseteq S$. They do not share opposite corners.
Hence they share $0$, $1$, or $2$ corners, and $2$ shared corners lie in one row
or in one column. We may assume that it is a row $a$. If they share no corner,
then $R_1$ is hidden, which is a $4$-cycle. If they share $(x,a)$ and $(y,a)$,
their other rows $b\ne c$ give the $4$-cycle $r_b\,c_x\,r_c\,c_y$. If they share
just $(x,a)$, write $R_1=\{x,y\}\times\{a,b\}$ and $R_2=\{x,z\}\times\{a,c\}$
with $y\ne z$ and $b\ne c$. The other $6$ corners are black and form the
$6$-cycle $r_a\,c_y\,r_b\,c_x\,r_c\,c_z$.
\end{proof}

The next lemma tells when a set of at most $2$ points is the presentation of
an axis-aligned rectangle.

\begin{lemma}[axis-aligned rivals]\label{lem:rival}
Let $S$ hide no axis-aligned rectangle, and let $P$ be a set of at most $2$
white points. A \emph{rival} of $P$ is an axis-aligned rectangle with
presentation $P$. The empty set has no rival. Otherwise $P$ has a rival exactly
in the following cases.
\begin{enumerate}[(a)]
\item $P=\{(x,r)\}$, and some $s\in B_x$ has $\{r,s\}\in F$.
\item $P=\{(x,r),(x,s)\}$ is a vertical pair, and $\{r,s\}\in F$.
\item $P=\{(x,r),(z,r)\}$ is a horizontal pair, and $B_x\cap B_z\ne\emptyset$.
\item $P=\{(x,r),(z,s)\}$ is an oblique pair with $x\ne z$ and $r\ne s$, and
$(x,s)$ and $(z,r)$ are black.
\end{enumerate}
\end{lemma}

\begin{proof}
A rival of $P\ne\emptyset$ has rows $\{r,s\}$ and columns $\{x,z\}$, and each of
its corners outside $P$ is black. Reading off these corners gives the $4$
conditions. In (a) the black corners are $(x,s)$, $(z,r)$, and $(z,s)$, so
$s\in B_x$ and column $z$ contains the pair $\{r,s\}$. Here $z\ne x$ holds
automatically, since $(x,r)$ is white.
\end{proof}

\subsection{What we use from the companion paper}

\begin{theorem}[{\cite{pemmasaniExtremal}}]\label{thm:imports}
Let $m\ge3$.
\begin{enumerate}[(a)]
\item \textup{(Row bound.)} Every axis-valid blackout of the $h\times m$ grid has
at most $m+q_h$ points. Hence $f(h,m)\le m+q_h$.
\item \textup{(Equality.)} Let $h\ge4$. A blackout of the $h\times m$ grid with
$m+q_h$ points is axis-valid if and only if every column is a singleton or a
double, the $q_h$ doubles have distinct row pairs, and these pairs are the edges
of a complete bipartite graph on the rows with parts of sizes $\lfloor h/2\rfloor$
and $\lceil h/2\rceil$.
\item \textup{(Values.)} We have $f(2,m)=m+1$ and $f(3,m)=m+2$. We have
$f(4,4)=7$, $f(4,5)=8$, and $f(4,m)=m+4$ for $m\ge6$. We have $f(5,5)=10$,
$f(5,6)=11$, $f(5,7)=12$, and $f(5,m)=m+6$ for $m\ge8$.
\end{enumerate}
\end{theorem}

In the companion paper the values $f(4,4)$, $f(4,5)$, $f(5,6)$, and $f(5,7)$
rest on computer searches, and $f(5,m)=m+6$ rests on a finite check that is
also verified in Lean. We do not need (c) for $2$ and $3$ rows, since
Theorem~\ref{thm:strip} and Corollary~\ref{cor:threevalue} prove these cases
again. Table~\ref{tab:values} lists $f(h,m)$ and $a_h(m)$ for $h\le5$.

\begin{table}[!htb]
\centering\footnotesize\setlength{\tabcolsep}{4pt}
\begin{tabular}{@{}r|rrrr@{}}
\multicolumn{5}{c}{$f(h,m)$}\\
\toprule
$m\backslash h$ & $2$ & $3$ & $4$ & $5$\\
\midrule
$2$ & $4$ & $4$ & $5$ & $6$\\
$3$ & $4$ & $5$ & $6$ & $7$\\
$4$ & $5$ & $6$ & $7$ & $8$\\
$5$ & $6$ & $7$ & $8$ & $10$\\
$6$ & $7$ & $8$ & $10$ & $11$\\
$7$ & $8$ & $9$ & $11$ & $12$\\
$8$ & $9$ & $10$ & $12$ & $14$\\
$9$ & $10$ & $11$ & $13$ & $15$\\
$10$ & $11$ & $12$ & $14$ & $16$\\
$11$ & $12$ & $13$ & $15$ & $17$\\
$12$ & $13$ & $14$ & $16$ & $18$\\
$13$ & $14$ & $15$ & $17$ & $19$\\
\bottomrule
\end{tabular}\hspace{1.2em}
\begin{tabular}{@{}r|rrrr@{}}
\multicolumn{5}{c}{$a_h(m)$}\\
\toprule
$m\backslash h$ & $2$ & $3$ & $4$ & $5$\\
\midrule
$2$ & $1$ & $12$ & $32$ & $80$\\
$3$ & $12$ & $53$ & $224$ & $892$\\
$4$ & $32$ & $224$ & $316$ & $1810$\\
$5$ & $80$ & $892$ & $1810$ & $8$\\
$6$ & $192$ & $3420$ & $26$ & $76$\\
$7$ & $448$ & $12704$ & $292$ & $1184$\\
$8$ & $1024$ & $45864$ & $2118$ & $4$\\
$9$ & $2304$ & $161992$ & $13332$ & $44$\\
$10$ & $5120$ & $562912$ & $75398$ & $350$\\
$11$ & $11264$ & $1930312$ & $402656$ & $3004$\\
$12$ & $24576$ & $6545496$ & $2037848$ & $24706$\\
$13$ & $53248$ & $21988784$ & $9894016$ & $193406$\\
\bottomrule
\end{tabular}
\caption{The values of $f(h,m)$ (\seqnum{A397315}) and $a_h(m)$
(\seqnum{A399580}) for $h\le5$. The columns $h=2$, $3$, $4$ of the counts are
\seqnum{A001787} from $m=3$, \seqnum{A401223}, and \seqnum{A401224}, and the
column $h=5$ is \seqnum{A401225} from $m=8$. The counts $a_4(m)$ and $a_5(m)$ with
$m\ge4$ are computer-assisted.}
\label{tab:values}
\end{table}
